% SPDX-FileCopyrightText: © 2026 Idiap Research Institute <contact@idiap.ch>
%
% SPDX-FileContributor: Andre Freitas <andre.freitas@idiap.ch>
% Neuro-symbolic AI Group
%
% SPDX-FileContributor: Danilo Gusicuma <danilo.gusicuma@idiap.ch>
% Neuro-symbolic AI Group
%
% SPDX-License-Identifier: MIT
% Generated from step-explanations.json
\paragraph{Step-by-step reconstruction.}
A proposed domination argument treats a small joint event probability as a history-uniform conditional ceiling. Test the finite law assigning probability one half each to outcomes 10 and 01.
The following is a source-guided illustrative reconstruction. Internal procedures are prescriptions; the native declarations do not establish their execution.
\begin{enumerate}[leftmargin=*,label=\arabic*.]
\item \textbf{Construct the conditional derivation} (\texttt{conditional\_\allowbreak{}derivation}).
\begin{description}[style=standard,labelindent=0pt,leftmargin=5.5em,labelwidth=5em,labelsep=0.5em,align=left,font=\normalfont\bfseries,itemsep=0.15em,parsep=0pt,topsep=0.3em]
\item[Purpose] Specify the mechanism’s required inequality and assumptions.
\item[Input] Bounded frontier; target; claim ceiling
\item[Method] The native producer is unresolved. Linked domination methods specify exposure history, retained boundary information and normalized conditional bounds; finite diagnostics attack the required implication.
\item[Check] Joint moments cannot replace history-uniform bounds; normalization and positive-probability conditioning must be justified.
\item[Output] Intended derivation packet: failed inference; open replacement obligation
\item[Next] Readiness assesses the packet; the declared projection is not a proof-producing algorithm.
\item[Example] For P(10)=P(01)=1/2, the two-site intersection probability is zero, but P(X2=1 given X1=0)=1. Thus the suggested conditional ceiling does not follow.
\end{description}
\item \textbf{Assess readiness for audit} (\texttt{derivation\_\allowbreak{}readiness}).
\begin{description}[style=standard,labelindent=0pt,leftmargin=5.5em,labelwidth=5em,labelsep=0.5em,align=left,font=\normalfont\bfseries,itemsep=0.15em,parsep=0pt,topsep=0.3em]
\item[Purpose] Assess whether the derivation is auditable.
\item[Input] Conditional packet; open-obligation reference
\item[Method] Check that the proposed inequality, retained state, assumptions and failure cases remain explicit rather than hidden by a certificate label.
\item[Check] A finite check cannot discharge full-measure, repeated-iteration or continuum hypotheses.
\item[Output] Derivation disposition: ready for scoped audit or quarantined
\item[Next] Original policy sends accepted packets to audit; native return omits the packet adapter.
\item[Example] Send the counterexample and the requirement for a normalized history-uniform bound to audit; withhold any domination conclusion inferred only from the zero intersection moment.
\end{description}
\end{enumerate}
\paragraph{Changed-input contrast.}
A proved normalized conditional bound for every admissible positive-probability exposure history would answer this particular objection. Small unconditional intersection moments alone do not, and even that bound would not settle continuum reconstruction.
